\documentclass[12pt, a4paper]{report}
\usepackage[T1]{fontenc}
\usepackage[utf8]{inputenc}
\usepackage[polish]{babel}
\usepackage{graphicx} % Required for inserting images
\usepackage[superscript,biblabel]{cite}

\usepackage{setspace}
\onehalfspacing
\usepackage[a4paper, margin=2.5cm]{geometry}

\title{Wieloklasowa kategoryzacja terenu na zdjęciach satelitarnych pod kątem przydatności do awaryjnego lądowania z uwzględnieniem specyfikacji technicznej wybranych statków powietrznych.}
\author{Jan Juszczyński}
\date{Maj 2026}

\begin{document}

\maketitle
\tableofcontents
\chapter{Wstęp}

\par Problematyka bezpieczeństwa lotniczego od lat stanowi temat obszernych dyskusji i polemik. Dopóki czynnik ludzki ma wpływ na przebieg lotu, dopóty wyeliminowanie nieobliczalności podejmowanych decyzji nie będzie możliwe. Warto jednak zauważyć, że algorytmy również nie są nieomylne, a samolot jest na tyle skomplikowaną maszyną, że potrzebny jest ktoś, kto będzie mógł zareagować w każdej, nawet najtrudniejszej sytuacji. 
\par Jednym z najważniejszych i najbardziej krytycznych momentów jest awaryjne lądowanie. Pilot ma wtedy zazwyczaj od kilku do kilkudziesięciu sekund na odpowiednią reakcję. Na podstawie położenia maszyny względem ziemi, etapu lotu, ilości paliwa w zbiornikach, przyczyn awarii i wielu innych aspektów, człowiek musi podjąć decyzję, od której często zależy życie jego oraz pasażerów, za których jest odpowiedzialny podczas lotu. 
\par Przebodźcowany komunikatami i sygnałami pilot może zareagować nieoptymalnie lub nawet błędnie. Sztuczna inteligencja może w ułamku sekundy zasugerować odpowiednie miejsce do lądowania, przyspieszając czas reakcji pilota i zwiększając szanse powodzenia operacji awaryjnego lądowania. Model, analizując zdjęcia satelitarne, jest w stanie rozpoznać potencjalne miejsce do awaryjnego lądowania samolotu. Architektura modelu pozwala nawet na stosunkowo proste zmodyfikowanie procesu uczenia tak, żeby sugestie bazowały na danym typie maszyny. Specyfikacja oraz masa operacyjna statku powietrznego (np. lekkiego samolotu turystycznego w kontrze do odrzutowca) determinują warunki niezbędne do bezpiecznego przyziemienia.
\par Liczne publikacje naukowe wskazują na wysoką skuteczność sieci neuronowych w klasyfikacji obszarów na rzecz rolnictwa, budownictwa, czy zagospodarowania terenu. Jednakże w temacie badań dotyczących zastosowania takiego modelu rozpoznawania terenu w lotnictwie i sytuacjach awaryjnych, istnieje poważna luka. 
\par W niniejszej pracy podjęto próbę przybliżenia tego tematu i udowodnienia, że sztuczna inteligencja jest w stanie pomóc pilotowi w skutecznym sprowadzeniu samolotu na ziemię i co za tym idzie - uratowaniu życia ludzi na pokładzie samolotu.

\chapter{Przegląd literatury i podstawy teoretyczne}


\section{Dobór terenu do wykonania procedury awaryjnego lądowania}
\par Odpowiedni dobór terenu pod awaryjne lądowanie jest kluczowym aspektem odpowiadającym za jego powodzenie. Z uwagi na brak możliwości weryfikacji wszystkich wariantów takich procedur w rzeczywistości, wskazane jest teoretyczne zaznajomienie pilota z najlepszymi lokalizacjami do wylądowania. Według wytycznych Federal Aviation Administration (FAA)\cite{faa_handbook}, niemal każdy teren może posłużyć jako odpowiednie lądowisko awaryjne, pod warunkiem, że pilot potrafi wykorzystać elementy maszyny do ochrony kabiny pasażerskiej i wytracenia energii kinetycznej. 
\par Istnieją trzy główne typy lądowań awaryjnych:
\begin{itemize}
    \item \textbf{Wymuszone lądowanie awaryjne:} Jest ono najbardziej niebezpieczne. Pilot zmuszony jest do podjęcia szybkiej decyzji np. z powodu pożaru albo niedziałającego silnika.

    \item \textbf{Lądowanie prewencyjne:}  Często stosowane w momencie, kiedy nie istnieje bezpośrednie ryzyko utraty kontroli nad maszyną, np. w momencie gdy instrumenty pokładowe kapitana i drugiego pilota będą wskazywały inne dane. Warto jednak zauważyć, że tego typu zdarzenie może się przeobrazić w wymuszone lądowanie awaryjne, jeśli pilot zbagatelizuje usterkę lub umyślnie będzie próbował ignorować ostrzeżenia samolotu.

    \item \textbf{Lądowanie na wodzie ("Ditching"):} Jest używane w ostateczności, kiedy nie ma szans na bezpieczny powrót na lotnisko lub samolot leci nad terenem gęsto zaludnionym i inny rodzaj lądowania spowodowałby zagrożenie zdrowia lub życia osób postronnych na ziemi. Manewr ten skutkuje głębokim uszkodzeniem maszyny. W ekstremalnych przypadkach również zatonięciem lub rozbiciem się o taflę wody.
\end{itemize}
\par Nie można pominąć faktu, iż dobór odpowiedniego miejsca do lądowania jest ściśle skorelowany z wielkością i specyfikacją maszyny. Siła przyziemienia i odległość do zatrzymania rosną proporcjonalnie do kwadratu prędkości. Z tego powodu lżejsza maszyna jest w stanie wylądować na krótkim kawałku łąki, podczas gdy większy samolot pasażerski potrzebuje większego fragmentu autostrady lub szerokiej rzeki.\cite{ntsb_hudson,ntsb_southern}
\par Poprawnie wykonane lądowanie awaryjne w terenie może być lepszą decyzją niż próba sprowadzenia samolotu na pas startowy. Jednakże w razie błędnej oceny może doprowadzić do katastrofy o jeszcze większej skali. W 1977 roku piloci linii Southern Airways, z powodu dużej wysokości od ziemi i presji, błędnie ocenili szerokość lokalnej autostrady i ruch na niej. W efekcie rozbili maszynę o stację benzynową, powodując jeszcze większą tragedię. \cite{ntsb_southern} Pilot, posiadając możliwość weryfikacji własnej oceny ze statystykami czy analizą satelitarną wygenerowaną przez sztuczną inteligencję, mógłby bardziej racjonalnie podejść do zaistniałej sytuacji i zwiększyć szanse powodzenia operacji awaryjnego lądowania.


% \section{Zastosowanie Konwolucyjnych Sieci Neuronowych (CNN) w klasyfikacji i analizie obrazów satelitarnych}

\section{Ewolucja metod klasyfikacji i analizy obrazów satelitarnych}

\par Przez ostatnie 50 lat analiza obrazów satelitarnych uległa znaczącej zmianie.\cite{cheng_remote_sensing} Obecnie możemy wyróżnić trzy główne kategorie klasyfikacji zdjęć na podstawie używanych metod:
\begin{itemize}
\item \textbf{Podejście klasyczne \emph{(Handcrafted Feature Based Methods)}:} jest to najstarszy sposób klasyfikacji obrazów. Polega na ręcznym projektowaniu filtrów i deskryptorów. Wymaga specjalistycznej wiedzy i dużej ilości czasu. Istnieje możliwość połączenia wielu metod w celu osiągnięcia lepszego rezultatu, jednakże przy bardziej złożonych danych łączenie metod staje się złożone i prowadzi do nieoptymalnych rezultatów.

\item \textbf{Uczenie bez nadzoru \emph{(Unsupervised Feature Learning Based Methods)}:} jest to pierwszy krok w stronę automatycznego wyodrębniania cech z surowych danych w celu usprawnienia procesu i standaryzacji wyników. Jednym z pierwszych napisanych algorytmów była Analiza Głównych Składowych (PCA). Metoda ta polega na redukcji danych do mniejszej liczby nieskorelowanych cech o największej możliwej wariancji. Nowszym algorytmem uczenia bez nadzoru jest algorytm k-średnich (k-means clustering). Polega on na gromadzeniu cech w klastrach, tak żeby wszystkie punkty w jednym klastrze były jak najbardziej podobne, a w innych - jak najbardziej różne. W porównaniu do podejścia klasycznego, uczenie bez nadzoru osiąga zdecydowanie lepsze wyniki, jednak z powodu braku użycia etykiet, metody te są ograniczone i mogą generować nieprawidłowe wyniki. 

\item \textbf{Uczenie głębokie \emph{(Deep Feature Learning Based Methods)}:} w 2006 roku doszło do przełomu w klasyfikacji obrazów. Zostały zaprojektowane pierwsze algorytmy uczenia głębokiego. Od tego czasu, naukowcy starają się usprawniać i optymalizować konwolucyjne sieci neuronowe (CNN), które pozwalają na najsprawniejszą i najbardziej zaawansowaną analizę obrazów przekształconych w tensory dzięki użyciu wielu warstw sieci neuronowych. Pierwsze warstwy zazwyczaj wykrywają kontury, podczas gdy kolejne rozpoznają kształty, wzory czy nawet kolory. Największą zaletą tego rozwiązania jest automatyzacja nauki i klasyfikacji zdjęć oraz użycie wielu warstw do analizy zdjęcia. Warto jednak zauważyć, że w porównaniu do uczenia bez nadzoru, uczenie głębokie jest procesem znacznie rozbudowanym, wymagającym większych zasobów do przetwarzania obrazów.
\end{itemize}
Wdrożenie modelu CNN w systemie lotniczym pozwala na szybkie przystosowanie algorytmu do lokalnej architektury poprzez douczanie (fine-tuning).

\section{Przegląd dotychczasowych prac badawczych i zbiorów danych}

\par Stale rosnące zaawansowanie sprzętu i algorytmów analizy obrazów pozwala badaczom i inżynierom na tworzenie i zastosowanie bardziej złożonych zbiorów danych. W przeszłości badacze mieli do dyspozycji bazy obrazów takie jak \emph{UC Merced Land-Use Dataset} i \emph{WHU-RS19 Dataset}.\cite{cheng_remote_sensing} Zawierały one 19 do 21 klas. Liczba taka pozwala na duże zróżnicowanie, lecz niestety do każdej klasy przyporządkowane było po 50 do 100 obrazów. Zbiory te z powodu niskiej rozdzielczości i małej liczby zdjęć na klasę (50-100) były całkowicie nieodpowiednie dla algorytmów uczenia maszynowego, prowadząc do zjawiska przeuczenia (overfitting).
\par Z czasem zaczęły powstawać specjalistyczne zbiory danych takie jak \emph{Brazilian Coffee Scene Dataset}. Zbiory tego typu są ukierunkowane na rozwiązanie konkretnego problemu. W tym przypadku znaczenie miało szczegółowe rozpoznanie czy obszar jest uprawą kawy czy nie. Obrazy są podzielone na dwie klasy w zależności od pokrycia obszaru polem kawowym. Jeśli zdjęcie zawierało ponad 85\% pola kawowego, było przyporządkowywane do klasy "kawa". Zbiór ten ma konkretne zastosowanie w rolnictwie i uprawie kawy. Wymagania systemów bezpieczeństwa lotniczego wymuszają zastosowanie zbiorów o znacznie większej różnorodności klas, obejmujących zarówno zróżnicowaną infrastrukturę jak i różne typy naturalnego terenu.
\par Odpowiednie i efektywne uczenie maszynowe wymaga wystarczająco dużego, szczegółowego i wieloklasowego zbioru danych. W efekcie w 2017 powstał \emph{EuroSat}.\cite{helber_euroSat} Zbiór ten jest jednym z największych, najbardziej różnorodnych benchmarków dostępnych do analizy zdjęć. Powstał on dzięki programowi Copernicus finansowanemu przez \emph{European Space Agency} (ESA). Składa się z 27.000 zdjęć podzielonych na 10 klas. Zdjęcia pochodzą z satelity \emph{Sentinel-2} umożliwiającej wielospektralne obrazowanie. Rejestracja danych w 13 różnych pasmach fal umożliwia dogłębniejszą analizę terenu na zdjęciach i otwiera nowe możliwości w uczeniu maszynowym pod kątem szukania miejsca do lądowania.

\chapter{Metodologia i przygotowanie zbioru danych}
\chapter{Architektura modelu głębokiego uczenia}
\chapter{Wyniki eksperymentu i podsumowanie}
\chapter{Wnioski}

\begin{thebibliography}{9}

\bibitem{faa_handbook}
  Federal Aviation Administration (FAA),
  \emph{Airplane Flying Handbook (FAA-H-8083-3C)}, U.S. Department of Transportation, Washington D.C., 2021.
\bibitem{ntsb_hudson}
  National Transportation Safety Board (NTSB),
  \emph{Aircraft Accident Report NTSB/AAR-10/03: Loss of Thrust in Both Engines After Encountering a Flock of Birds and Subsequent Ditching on the Hudson River US Airways Flight 1549}, Washington D.C., 2010.
\bibitem{ntsb_southern}
    National Transportation Safety Board (NTSB),
    \emph{Aircraft Accident Report NTSB-AAR-78-3: Southern Airways Inc., DC-9-31, N1335U, New Hope, Georgia, April 4, 1977}, Washington D.C., 1978.
\bibitem{cheng_remote_sensing}
    G. Cheng, J. Han, and X. Lu,
    \emph{Remote Sensing Image Scene Classification: 
    Benchmark and State of the Art}, Proceedings of the IEEE, Volume: 105, Issue: 10, October 2017, pp. 1865--1883, DOI: 10.1109/JPROC.2017.2675998.
\bibitem{helber_euroSat}
    P. Helber, B. Bischke, A. Dengel, and D. Borth,
    \emph{EuroSAT: A Novel Dataset and Deep Learning
Benchmark for Land Use and Land Cover
Classification}, IEEE Journal of Selected Topics in Applied Earth Observations and Remote Sensing, Volume: 12, Issue: 7, July 2019, pp. 2217--2226, DOI: 10.1109/JSTARS.2019.2918242.


\end{thebibliography}

\end{document}
